% TOP_PROBLEM: {"id": 30002360, "title": "Grothendieck–Serre Conjecture for Regular Local Rings", "category_id": 4, "category": {"id": 4, "name": "algebra", "display_name": "Algebra", "description": "Group theory, ring theory, field theory, and algebraic structures.", "slug": "algebra", "order_index": 4, "created_at": "2026-07-31T15:26:25.671Z"}, "status": "partially_solved"}
% FIRST_POSTED: 2026-09-25
% ENTRY_KIND: statement_only
% !TeX program = lualatex
% Definition and sources only; this file is not an LLM solution attempt.
\documentclass[11pt]{article}
\usepackage[margin=25mm]{geometry}
\usepackage{fontspec}
\setmainfont{Latin Modern Roman}
\usepackage{amsmath,amssymb,mathtools}
\usepackage{unicode-math}
\setmathfont{Latin Modern Math}
\usepackage{xurl,hyperref}
\hypersetup{colorlinks=true,urlcolor=blue}
\setcounter{secnumdepth}{0}
\setlength{\parindent}{0pt}
\setlength{\parskip}{5pt}
\setlength{\emergencystretch}{3em}
\begin{document}
\section*{\texorpdfstring{Grothendieck–Serre Conjecture for Regular Local Rings}{Grothendieck–Serre Conjecture for Regular Local Rings}}\label{30002360}
\subsection{Definitions and mathematical statement}
A Noetherian local ring $(R,\mathfrak m)$ is regular when
$\dim R=\dim_{R/\mathfrak m}\mathfrak m/\mathfrak m^2$; it is a domain here, with fraction field $K$. A reductive $R$-group scheme is smooth affine with connected reductive geometric fibers. A $G$-torsor is a principal homogeneous space locally trivial for the \"etale topology. The assertion is
\[
 \ker\bigl(H^1_{\mathrm{et}}(R,G)\to H^1_{\mathrm{et}}(K,G)\bigr)=\{1\},
\]
where these are pointed sets and the kernel is the preimage of the trivial torsor. Thus generic triviality forces triviality over $R$. The exact statement does not impose a field inside $R$, excellence, isotropy, or a constant group scheme.
\par\smallskip\textit{Formulation and concept references:} \href{https://arxiv.org/pdf/2201.06424v4}{[S1]}.

\subsection{Short English statement}
If a principal bundle under a reductive group over a regular local ring is trivial over the fraction field, must it already be trivial over the ring?
\subsection{Sources}
\begin{itemize}
\item[S1]Problems about torsors over regular rings, with an appendix by Yifei Zhao; arXiv2201.06424v4,5May2025; HTML internal29March2023.
\url{https://arxiv.org/pdf/2201.06424v4}.
\textit{Formulation source.}
\end{itemize}
\end{document}
